% =====================================================================================
% Chapter 11 · Threshold perks — regression discontinuity
%
% Built from notebooks/09_threshold_perk_rdd.ipynb. NOT ONE NUMBER IN THIS FILE IS TYPED
% BY HAND: every figure in the prose is a macro from build/macros.tex, generated by
% cmp.report from the executed notebook. The only literals below are (i) the constants of
% the data-generating model in §11.2, which are *definitions* rather than results, and
% (ii) the decision-rule thresholds (0.9 / 0.5) and the |t| < 1.645 falsification bar,
% which are *conventions* stated in the text, not measurements.
% =====================================================================================
\chapter{Threshold perks: regression discontinuity}
\label{ch:rdd}

\begin{quote}\small
\emph{Cross \eur{\nbNineCutoff} of annual spend and the loyalty programme grants Gold status, whose
perks cost the firm \eur{\nbNinePerkCost} per member against a member-year worth
\eur{\nbNineAnnualValue}. The dashboard says Gold members retain \nbNineNaive{}\,pp better than
everyone else, with an interval of $[\nbNineNaiveLo, \nbNineNaiveHi]$ so tight it will pass every
review --- and the planted truth, \nbNineTrueJump{}\,pp, lies \nbNineNaiveVerdict{} it. Nothing in a
confidence interval knows about confounding. Regression discontinuity throws away
\nbNineRdDropped{} of the \nbNineN{} customers, keeps only those within \eur{\nbNineH} of the
threshold, and returns \nbNineRd{}\,pp $[\nbNineRdLo, \nbNineRdHi]$, which contains the truth ---
and it survives a density test, two covariate-continuity tests, \nbNinePlaceboN{} placebo cutoffs
and a polynomial-order check. It does not survive the bandwidth. Across \nbNineBwNWindows{}
windows the estimate ranges \nbNineBwEstLo{} to \nbNineBwEstHi{}\,pp and the significance verdict
itself flips, \nbNineBwNInsignif{} of \nbNineBwNWindows{} windows being unable to exclude zero. The
union of those intervals, $[\nbNineBwUnionLo, \nbNineBwUnionHi]$, is
\nbNineBwUnionOverWidth$\times$ the width of the one interval the analysis reports.
\Cref{sec:rd:wrong} is about the number that interval is not showing you.}
\end{quote}

% =====================================================================================
\section{The decision}
\label{sec:rd:decision}

A loyalty programme grants \emph{Gold} status to any customer whose annual spend crosses
\eur{\nbNineCutoff}. Gold carries perks --- free delivery, an early-access window, a service line
--- which cost the firm \eur{\nbNinePerkCost} per member per year. The retention team's dashboard
reports that Gold members retain \nbNineNaive{}\,pp better than everyone else. The proposal on the
table is to keep the programme, and the questions the CMO actually asks are two: \textbf{does the
perk cause retention, or do big spenders simply retain anyway?} And \textbf{if it does, is
\eur{\nbNinePerkCost} the right price for it?}

The cost of a wrong answer is not symmetric with the cost of a right one, and it is not small. A
perk that does nothing is a permanent line in the P\&L, renewed every year on the strength of a
number that was never a causal number. A perk that works and is killed hands a retention lever to a
competitor. And a third failure --- the one this chapter is really about --- is more insidious than
either: a perk that works, priced wrongly, defended by an interval that looked authoritative.

The naive comparison is the problem in miniature, and it is worth being precise about \emph{why} it
fails rather than merely asserting that it does. Gold members are, by the definition of the tier,
the firm's biggest spenders. Big spenders retain better for reasons that have nothing to do with
perks: they are more committed, longer tenured, more habituated. So the Gold-vs-rest gap is the
perk's effect \emph{plus} the whole climb of retention against spend --- the confounder --- and the
two are not separable by staring at the number. \Cref{fig:rd:confound} shows the climb; the gap
between the two dashed group means is what the dashboard reports.

\input{build/floats/nb09_confound}

The naive estimate is \nbNineNaive{}\,pp with a 90\,\% interval of
$[\nbNineNaiveLo, \nbNineNaiveHi]$ --- a standard error of \nbNineNaiveSe{}\,pp on
\nbNineN{} customers, computed with an honest \nbNineNaiveCov{} covariance. It is
\emph{arithmetically impeccable and causally worthless}. The planted truth, \nbNineTrueJump{}\,pp,
falls \nbNineNaiveVerdict{} that interval, and the estimate is inflated by \nbNineNaiveBias{}\,pp
of pure selection.

Fix the reader's attention on the shape of that failure, because it recurs throughout applied
marketing analytics: \textbf{the interval is tight, and the tightness is exactly what makes it
dangerous}. A wide interval around a wrong number invites scrutiny. A narrow one does not. And
nothing in the construction of a confidence interval consults the causal structure of the problem:
an interval prices \emph{sampling} noise, on the assumption that the estimator is aimed at the
estimand. Aim it at the wrong estimand and the interval will shrink faithfully around the wrong
number as $n$ grows. Confounding is not a variance problem, and no amount of data fixes it.

What is needed instead is a comparison in which the confounder does not vary. That is what a sharp
threshold hands us for free.

% =====================================================================================
\section{The data-generating model}
\label{sec:rd:dgm}

A causal method cannot be graded on real data, because on real data the counterfactual is missing
--- that is the entire problem. So the method is first graded against a world whose answer is known
by construction.

The simulator draws \nbNineN{} customers. Each has an annual spend $S_i$ --- the \emph{running
variable}, also called the forcing or assignment variable --- concentrated around the cutoff
$c = \nbNineCutoff$ so that both sides of it are densely populated:
%
\begin{equation}
  S_i \;\sim\; \mathcal N\!\left(100,\ 40^{2}\right), \ \text{clipped below at } 1,
  \qquad
  D_i \;=\; \mathbf 1\!\left[\,S_i \ge c\,\right].
  \label{eq:rd:running}
\end{equation}
%
$D_i$ is the Gold indicator, and \cref{eq:rd:running} says the assignment rule is \emph{sharp}:
crossing $c$ grants the perk with probability one, and not crossing it grants the perk with
probability zero. There is no discretion, no application form, no partial compliance. (A programme
in which crossing the threshold merely raises the \emph{probability} of the perk is a \emph{fuzzy}
design; \cref{sec:rd:ident} says what changes.)

Retention is Bernoulli, with a probability that rises \emph{smoothly and linearly} in spend and
jumps \emph{discontinuously} at the cutoff:
%
\begin{align}
  p(S_i, D_i) &\;=\; \operatorname{clip}\!\Big(
      \underbrace{0.45}_{\text{baseline}}
      \;+\; \underbrace{0.0015\,(S_i - c)}_{\text{the confounder}}
      \;+\; \underbrace{\tau\, D_i}_{\text{the perk}} ,\ 0.02,\ 0.98 \Big),
      \label{eq:rd:p}\\[2pt]
  Y_i \mid S_i, D_i &\;\sim\; \text{Bernoulli}\big(p(S_i, D_i)\big),
  \qquad \tau \;=\; 0.14 .
      \label{eq:rd:y}
\end{align}
%
Two terms, and they are the whole chapter. The slope $0.0015$ per euro --- \pc{0.15} of retention
per euro of spend --- is the \emph{smooth} part: big spenders retain more anyway, and this term
knows nothing about any perk. The jump $\tau = \nbNineTrueJump{}\,\text{pp}$ is the \emph{causal}
part, and it is discontinuous by construction. The planted \nbNineTrueJump{}\,pp is the grade every
estimator in this chapter is marked against, and it is the last time in the chapter that anything
is known for certain.

Two properties of \crefrange{eq:rd:running}{eq:rd:y} deserve to be named out loud, because on real
data they are assumptions rather than facts, and \cref{sec:rd:ident} will have to assume them.
First, \emph{the smooth part really is smooth at $c$}: the confounder passes through the threshold
without a kink or a break, so any break in observed retention at $c$ can only be the perk. Second,
\emph{nobody manipulates $S$}: spend is drawn before assignment and no customer nudges a basket to
clear the tier, so the density of \cref{eq:rd:running} is continuous at $c$.

The confounder is what makes the naive comparison fail, and the size of that failure is available
in closed form here, which is a rare luxury. With $\Delta = S - c$ symmetric about zero and
$D = \mathbf 1[\Delta \ge 0]$,
%
\begin{equation}
  \underbrace{\E[Y \mid D{=}1] - \E[Y \mid D{=}0]}_{\text{the dashboard's number}}
  \;=\; \underbrace{\tau}_{\text{causal}}
  \;+\; \underbrace{0.0015\,\big(\E[\Delta \mid \Delta \ge 0] - \E[\Delta \mid \Delta < 0]\big)}
        _{\text{selection}} ,
  \label{eq:rd:decomp}
\end{equation}
%
and for a normal running variable the selection term evaluates to
$0.0015 \times 2\sigma\sqrt{2/\pi} = \nbNineSelectionClosedForm{}\,\text{pp}$. The measured
naive bias is \nbNineNaiveBias{}\,pp; the remainder is clipping and sampling noise. The dashboard is
not wrong by accident or by a little. It is wrong by a quantity that can be predicted before the
data are seen, and it nearly doubles the truth.

% =====================================================================================
\section{Identification}
\label{sec:rd:ident}

\subsection*{The estimand}

In the potential-outcome notation of \citet{rubin1974} and \citet{imbens2015}, write $Y_i(1)$ for
customer $i$'s retention \emph{with} the perk and $Y_i(0)$ for retention \emph{without} it. Exactly
one is ever observed, and which one is decided by $D_i$ in \cref{eq:rd:running}. The object this
chapter estimates is the effect of the perk \textbf{for customers standing exactly at the
threshold}:
%
\begin{equation}
  \tau_{\mathrm{RD}} \;=\; \E\big[\,Y(1) - Y(0) \,\big|\, S = c \,\big].
  \label{eq:rd:estimand}
\end{equation}
%
Read \cref{eq:rd:estimand} slowly, because the conditioning event is the entire method and also its
entire limitation. It is a \emph{local} average treatment effect --- local not to a subpopulation
defined by compliance, as in an instrumental-variables design, but local to a \emph{point on the
running variable} \citep{hahn2001, imbens2008}. It is the perk's effect on the customer who spends
\eur{\nbNineCutoff}. It is silent about the \eur{500} whale, and it is silent about the \eur{40}
casual. Every number in this chapter --- classical, Bayesian, both intervals, and the euro decision
--- targets the single scalar of \cref{eq:rd:estimand}, and keeping the estimand fixed is what makes
\cref{sec:rd:compare} a comparison rather than a change of subject.

\subsection*{The idea, in one sentence}

\Citet{thistlethwaite1960} saw it first, in a study of merit certificates awarded on a test-score
cutoff. \emph{Just below and just above a sharp threshold, the units are essentially identical.} A
customer who spent \eur{99} and a customer who spent \eur{101} are the same kind of person: the same
tenure, the same order history, the same latent propensity to stick around. Whatever determined
which side of the line they landed on --- a rounding of a basket, one more purchase in December, a
gift they happened to buy --- is, near the threshold, as good as a coin flip. The only thing that
changes discontinuously at \eur{\nbNineCutoff} is the perk. So compare them.

That intuition has a formal name --- \emph{local randomization} \citep{lee2010, cattaneo2016} ---
and it converts an observational study into something that behaves, in a neighbourhood of the
cutoff, like an experiment nobody ran.

\subsection*{The assumptions}

Identification of \cref{eq:rd:estimand} rests on the following. They are stated here, and each is
either tested in \cref{sec:rd:valid} or explicitly declared untestable --- an assumption with no
accompanying test is a wish. \Cref{tab:rd:assumptions} is the audit, and its rows A1--A5 are these
five, in this order.

\begin{assumption}[Continuity of the potential-outcome means]\label{as:rd:cont}
The functions $s \mapsto \E[Y(0) \mid S = s]$ and $s \mapsto \E[Y(1) \mid S = s]$ are continuous at
$s = c$ \citep{hahn2001}. In words: absent the perk, nothing about retention would have jumped at
\eur{\nbNineCutoff}. Under \cref{as:rd:cont},
%
\begin{equation}
  \tau_{\mathrm{RD}}
  \;=\; \lim_{s \downarrow c} \E[Y \mid S = s] \;-\; \lim_{s \uparrow c} \E[Y \mid S = s],
  \label{eq:rd:limits}
\end{equation}
%
and the right-hand side is made of observable quantities alone. That is the whole identification
result: \cref{eq:rd:limits} turns an unobservable contrast of potential outcomes into a difference
of two one-sided limits of a regression function. \textbf{\Cref{as:rd:cont} cannot be tested
directly} --- it is a statement about $Y(0)$ above the cutoff and $Y(1)$ below it, and neither is
ever observed. Only its observable implications can be tested, and \cref{sec:rd:valid} tests two of
them.
\end{assumption}

\begin{assumption}[No manipulation of the running variable]\label{as:rd:manip}
Customers cannot precisely sort themselves across $c$. If a customer who would have spent \eur{97}
can see the threshold and top up to \eur{101}, then the customers just above the line are the ones
who \emph{wanted} the perk, and wanting the perk is correlated with staying. \Cref{as:rd:cont} dies,
and it dies in the direction that flatters the programme. The observable implication is that the
\emph{density} of $S$ is continuous at $c$ --- a pile-up just above the threshold is the
fingerprint of sorting \citep{mccrary2008}. Tested in \cref{sec:rd:valid}.
\end{assumption}

\begin{assumption}[No other rule at the cutoff]\label{as:rd:only}
The perk is the \emph{only} thing that switches on at \eur{\nbNineCutoff}. If the same threshold
also triggers a different email cadence, a different service tier, or a different discount, then
\cref{eq:rd:limits} identifies the effect of the \emph{bundle} and calling it ``the perk'' is a
naming error with a budget attached. This is partly an institutional question --- go and read the
programme rules --- and partly testable: \cref{sec:rd:valid} refits the estimator at
\nbNinePlaceboN{} thresholds where nothing is supposed to happen, and finds nothing.
\end{assumption}

\begin{assumption}[The local polynomial is adequate]\label{as:rd:poly}
Within the estimation window, the conditional mean on each side of the cutoff is well approximated
by the fitted polynomial (here, a straight line). This is not an assumption about the world so much
as about the \emph{estimator}: if the true regression function curves sharply inside the window and
a line is fitted through it, the line's intercept at $c$ is biased and the bias is booked as
effect. Probed in \cref{sec:rd:valid} by refitting with a quadratic, and discussed --- because the
conventional remedy is worse than the disease --- in \cref{sec:rd:wrong}.
\end{assumption}

\begin{assumption}[The bandwidth is the right one]\label{as:rd:bw}
The estimate is read from customers within $\pm h$ of the cutoff, and the answer does not depend
materially on $h$. \textbf{This is the assumption this chapter cannot discharge}, and
\cref{sec:rd:wrong} is devoted to it. It is stated as an assumption, rather than buried as an
implementation detail, precisely because that is what it is.
\end{assumption}

\begin{remark}[Sharp and fuzzy]
\Cref{eq:rd:running} makes assignment deterministic. When the threshold instead shifts the
\emph{probability} of treatment --- customers above \eur{\nbNineCutoff} are \emph{offered} Gold and
some decline --- the design is \emph{fuzzy}, and the jump in the outcome must be divided by the
jump in take-up. The result is a Wald ratio, and the design becomes an instrumental-variables
problem in which the threshold indicator instruments for treatment
\citep{hahn2001, imbens1994}; the endogenous-exposure chapter treats that machinery in full. This
chapter's programme is sharp, so the denominator is one.
\end{remark}

% =====================================================================================
\section{The classical baseline}
\label{sec:rd:classical}

Before any sampler, the discipline of this book is to ask what a competent analyst produces
\emph{without} one. The answer is not a different analysis. It is the \emph{same} estimand of
\cref{eq:rd:estimand}, identified by the \emph{same} \cref{as:rd:cont}, estimated with the simplest
tool that is still correct. The causal work lives in the identification argument, never in the
machinery, and it is worth seeing that once, cleanly, before a likelihood arrives to obscure it.

\subsection*{The estimator}

\Cref{eq:rd:limits} asks for two one-sided limits of a regression function. Estimate each by a
\emph{local linear} regression: fit a straight line to the customers within $h$ euros below the
cutoff, another to those within $h$ euros above, and read the vertical gap where they meet.
Equivalently, and this is how it is fitted, run one weighted least-squares regression on the window
$|S_i - c| \le h$,
%
\begin{equation}
  Y_i \;=\; \alpha \;+\; \tau\, D_i \;+\; \beta\,(S_i - c) \;+\; \gamma\, D_i \,(S_i - c)
            \;+\; \varepsilon_i ,
  \label{eq:rd:ols}
\end{equation}
%
with weights $w_i = 1 - |S_i - c| / h$ --- a \emph{triangular kernel}, which gives full weight to a
customer standing on the threshold and zero weight to one at the window's edge. Because the running
variable is centred at $c$, the coefficient $\hat\tau$ \emph{is} the jump at the cutoff, and its
standard error is the standard error of a single regression coefficient. The interaction term
$\gamma$ earns its place: it lets the slope differ on the two sides, so that a change in the
\emph{trend} at the threshold is not mistaken for a change in the \emph{level}.

Three choices in \cref{eq:rd:ols} are worth defending, because each is a place where an RD analysis
can be quietly wrong.

\begin{itemize}[leftmargin=*, itemsep=.3em]
  \item \textbf{Local, not global.} A single polynomial fitted to \emph{all} \nbNineN{} customers
        would let a \eur{400} whale --- a customer with nothing whatever to say about the perk's
        effect at the margin --- pull the fitted curve at \eur{\nbNineCutoff}.
        \Cref{eq:rd:estimand} is a statement about $S = c$, so only customers near $c$ should have
        a vote. \Cref{sec:rd:wrong} shows what happens when that discipline is abandoned.
  \item \textbf{Linear, not high-order.} A straight line inside a narrow window is a first-order
        Taylor expansion of whatever the true regression function is, and the narrower the window,
        the better the approximation. Adding curvature is a remedy that costs more than the disease
        \citep{gelman2019}; \cref{sec:rd:wrong} returns to it.
  \item \textbf{Heteroskedasticity-robust standard errors.} The outcome is Bernoulli, so its
        variance is $p(1-p)$ and depends on $S$ mechanically. Homoskedastic OLS errors are
        therefore known to be wrong before the data are seen, and HC1 errors cost nothing.
        The reported covariance is \nbNineRdCov.
\end{itemize}

\subsection*{The result, and what it cost}

Fitted at $h = \eur{\nbNineH}$, \cref{eq:rd:ols} gives a jump of \textbf{\nbNineRd{}\,pp} with a
90\,\% confidence interval of $[\nbNineRdLo, \nbNineRdHi]$ (standard error \nbNineRdSe{}\,pp). The
planted \nbNineTrueJump{}\,pp falls \nbNineRdVerdict{} it; the point estimate is off by
\nbNineRdErr{}\,pp, which \cref{sec:rd:valid} shows to be sampling noise rather than bias.

Now the price. The window contains \nbNineRdN{} customers --- \nbNineRdNLeft{} below the cutoff and
\nbNineRdNRight{} above it --- which means the estimator \textbf{discarded \nbNineRdDropped{} of the
\nbNineN{} customers}, \pc{\nbNineRdDroppedPct} of the base, in order to buy comparability. It paid
for that in width: the RD interval is \nbNineWidthVsNaive$\times$ as wide as the naive one
(\cref{tab:rd:naive}).

\input{build/tables/nb09_naivetable}

\Cref{tab:rd:naive} is the chapter in one table, and the lesson generalises well beyond regression
discontinuity. \textbf{The naive estimator is more precise and it is wrong. The RD estimator is
less precise and it is right.} Precision and accuracy are different virtues, they trade against
each other here, and only one of them can be improved by collecting more data. An analyst who
optimises for a tight interval will choose the naive number every time, and will be able to defend
that choice with statistics all the way to the point at which the firm has spent a decade renewing
a perk on the strength of a selection effect.

\subsection*{What the classical arm cannot say}

One thing, and \cref{sec:rd:euro} needs it. A 90\,\% confidence interval does \emph{not} say that
there is a 90\,\% probability the true jump lies inside it. It is a property of the \emph{procedure}
--- across repeated samples, 90\,\% of intervals built this way cover the truth --- and this
particular interval either contains the truth or does not. The decision rule the CMO will apply,
$\Prob(\text{perk pays}) \ge 0.9$, is a probability statement \emph{about the effect}, and no
amount of frequentist effort produces one. That is the gap the next section exists to fill, and it
is the only gap it fills.

% =====================================================================================
\section{The Bayesian model}
\label{sec:rd:bayes}

The Bayesian arm re-houses \cref{eq:rd:ols} rather than replacing it. The design matrix is
identical --- the same window, the same centred running variable, the same treatment indicator and
the same interaction --- and what changes is that the coefficients acquire a distribution instead of
a point:
%
\begin{align}
  Y_i \mid \beta, \sigma
    &\;\sim\; \mathcal N\!\big(\mu_i,\ \sigma^{2}\big),
    \qquad |S_i - c| \le h, \label{eq:rd:lik}\\
  \mu_i &\;=\; \beta_0 \;+\; \beta_1 S_i \;+\; \beta_2 D_i \;+\; \beta_3 D_i S_i ,
    \label{eq:rd:mu}\\
  \beta &\;\sim\; \mathcal N(0,\ 50^2 \mathbf I), \qquad \sigma \;\sim\; \text{HalfNormal}(1),
    \label{eq:rd:prior}
\end{align}
%
and fitted with CausalPy's \code{RegressionDiscontinuity}, on a PyMC linear-regression
backend \citep{salvatier2016}. The jump is the posterior of the fitted discontinuity at $s = c$.

Three modelling choices need justifying, and one of them is the kind of thing that gets waved
through.

\textbf{The priors do nothing, and that is deliberate.} A normal prior of scale $50$ on
coefficients whose posteriors live in $[0, 1]$ is, for practical purposes, flat: it excludes nothing
a retention probability could plausibly be. This is the book's recurring finding stated once more
--- \emph{the prior is not where the action is}. With \nbNineNWindow{} customers in the window, the
likelihood \cref{eq:rd:lik} determines the answer, and \cref{sec:rd:compare} measures exactly how
little the Bayesian layer moved it.

\textbf{The likelihood is Gaussian on a binary outcome} --- a \emph{linear probability model}. That
looks like an error and is not, for a specific reason worth stating. The estimand
\cref{eq:rd:estimand} is a difference of two \emph{conditional means} at a single point, and a
conditional mean of a Bernoulli variable is a probability. Nothing in \cref{eq:rd:limits} requires a
link function; the link only matters if one wants the fitted curve to be sensible far from the data,
and the entire method refuses to look far from the data. The claim is checkable, so it was checked:
refitting the same window with a logistic link gives a jump of \nbNineLogitJump{}\,pp against the
linear model's \nbNineBayes{}\,pp, a difference of \nbNineLogitDiff{}\,pp. Inside a narrow window,
the link choice is immaterial. (The Gaussian likelihood \emph{is} heteroskedastic-misspecified ---
Bernoulli variance is $p(1-p)$, not constant $\sigma^2$ --- which is precisely what the classical
arm's HC1 errors correct for and the posterior does not. \Cref{sec:rd:compare} finds the resulting
distortion to be negligible here, and says why.)

\textbf{The sampler converged.} \nbNineNDraws{} post-warmup draws by the No-U-Turn sampler
\citep{hoffman2014}, with $\hat R = \nbNineRhat$, a minimum bulk effective sample size of
\nbNineEss{} \citep{vehtari2021}, and \nbNineDivergences{} divergent transitions. Nothing that
follows is a sampler artefact.

The posterior mean jump is \nbNineBayes{}\,pp with a 90\,\% credible interval of
$[\nbNineBayesLo, \nbNineBayesHi]$ --- posterior standard deviation \nbNineBayesSd{}\,pp.
\Cref{fig:rd:fit} is the picture the whole method lives in: two local-linear fits, one per side,
and \textbf{the height of the step where they meet at the cutoff is the estimate}.

\input{build/floats/nb09_rdfit}

Compare that picture with \cref{fig:rd:confound} and the difference between the two analyses is
visible rather than argued. The naive comparison contrasts the average \emph{height} of the orange
cloud with the average height of the green one, and credits the perk with the entire upward climb.
Regression discontinuity reads only the vertical \emph{break} at \eur{\nbNineCutoff}, so the climb
--- the confounder --- cancels. Same data, same customers, two numbers that differ by
\nbNineNaiveBias{}\,pp.

% =====================================================================================
\section{Diagnostics and validation}
\label{sec:rd:valid}

RD's diagnostics are unusually good, and this section is long because they deserve to be. The
reason they are good is structural: the design makes claims about \emph{observable} quantities ---
the density of the running variable, the behaviour of pre-determined covariates, the absence of
jumps where no rule exists --- and each of those claims is a test that could fail. Most causal
methods cannot say that. \Cref{tab:rd:assumptions} is the audit in full.

\input{build/tables/nb09_assumptions}

\subsection*{A2 --- is the running variable manipulated?}

This is the assumption that kills real RD studies, and it kills them silently. If customers can see
the threshold and cross it deliberately --- a \eur{97} basket topped up to \eur{101} to secure free
delivery --- then the customers just above the line are systematically the ones who \emph{wanted}
the perk. Wanting the perk is correlated with intending to stay. \Cref{as:rd:cont} fails, and the
RD estimate inherits exactly the selection bias it was built to escape.

The fingerprint is a discontinuity in the \emph{density} of $S$ at $c$: manipulators move from just
below to just above, so the mass just below is depleted and the mass just above piles up
\citep{mccrary2008}. \Cref{fig:rd:density} runs the test. The density just below the cutoff is
\nbNineMccBelow{} customers per euro and just above it is \nbNineMccAbove{}, a log-ratio of
\nbNineMccLogRatio{} and a binomial $z$ of \nbNineMccZ{} on a $\pm$\eur{\nbNineMccBw} window. No
pile-up, no evidence of sorting.

\input{build/floats/nb09_density}

Two honest qualifications. The test implemented here is a \emph{local-constant} binomial
approximation to \citeauthor{mccrary2008}'s test, not the full local-linear density estimator with
its Wald standard error; it is adequate as a ``did anyone pile up?'' screen and should be cited as
such. And it is a test with power against \emph{precise} manipulation only. A customer who nudges
spend upwards by a vague and noisy amount does not produce a visible discontinuity, and no density
test will find them. This is why \cref{as:rd:manip} is a claim about the business, checked against
the data, and not a claim the data can settle alone: know whether the threshold is visible to
customers, whether the tier is granted in real time or at year end, and whether anything about the
basket can be adjusted after the fact.

\subsection*{A1 --- do pre-determined covariates jump?}

\Cref{as:rd:cont} is untestable, but it has testable implications, and this is the sharpest of them.
Anything determined \emph{before} the perk was granted --- how long the customer has been with the
firm, how many orders they placed last year --- \emph{cannot} have been caused by it. So run the
identical RD machinery of \cref{eq:rd:ols} with a pre-determined covariate as the outcome. The
estimated jump at \eur{\nbNineCutoff} must be zero. If it is not, then the customers either side of
the threshold differ in something other than the perk, and the comparison the whole method rests on
is not a fair one.

This is the RD analogue of an experiment's balance table, and it is exactly as informative.
\Cref{fig:rd:covariates} shows both covariates climbing steeply with spend --- Gold members
\emph{are} longer-tenured and heavier orderers, which is precisely why the dashboard's number lied
--- and both passing smoothly through the cutoff. The jumps are \nbNineCovTenure{} months
($t = \nbNineCovTenureT$) and \nbNineCovOrders{} orders ($t = \nbNineCovOrdersT$): zero, to within
noise, against the conventional $|t| < 1.645$ bar.

\input{build/floats/nb09_covariates}

\subsection*{A3 --- placebo cutoffs, or: does the machinery find jumps that are not there?}

An estimator that reports a discontinuity at \eur{\nbNineCutoff} is only interesting if it
\emph{fails} to report one at \eur{70}. So the identical estimator is refit at \nbNinePlaceboN{}
fake thresholds, four below the real cutoff and four above, each in a $\pm$\eur{\nbNinePlaceboBw}
window, and each restricted to a single side of the real cutoff so that no placebo window can catch
sight of the genuine jump.

\input{build/tables/nb09_placebogrid}
\input{build/floats/nb09_placebo}

Nothing is found. Every one of the \nbNinePlaceboN{} placebo intervals covers zero, and the largest
$t$-statistic anywhere on the grid is \nbNinePlaceboMaxT{} --- at \eur{\nbNinePlaceboMaxCut}, where
the estimator reports \nbNinePlaceboMaxJump{}\,pp and cannot distinguish it from noise
(\cref{tab:rd:placebo}, \cref{fig:rd:placebo}). Only the real threshold clears zero.

The logic of this test is worth spelling out, because it is often run and rarely understood. A
placebo that \emph{did} find a jump would not merely be an anomaly to be explained away. It would
demonstrate that the estimator manufactures discontinuities out of curvature, noise or a badly
chosen window --- in which case the jump at \eur{\nbNineCutoff} would be evidence of nothing at all,
however large it was. The placebos are a test of the \emph{instrument}, not of the world.

\subsection*{A4 --- polynomial order and the functional form}

Refit \cref{eq:rd:ols} inside the same window with a quadratic on each side. The jump moves from
\nbNinePolyLinear{}\,pp to \nbNinePolyQuadratic{}\,pp --- a change of \nbNinePolyDiff{}\,pp, well
inside the standard error of either. The linear approximation is adequate at this bandwidth, which
is what one expects when the true conditional mean is a straight line (\cref{eq:rd:p}) and the
window is narrow. \Cref{sec:rd:wrong} takes up the question the other way around --- what happens
when an analyst reaches for a \emph{high}-order global polynomial instead --- and the answer is
worse than most practitioners expect.

\subsection*{Recovery: is the estimator unbiased, and is its interval calibrated?}

The single-sample estimate sits \nbNineRdErr{}\,pp below the truth. That is either bias, which would
be a defect of the estimator, or noise, which would not, and one sample cannot tell the difference.
So the whole procedure --- simulate, fit, interval --- is re-run on \nbNineNSeeds{} fresh samples
from \crefrange{eq:rd:running}{eq:rd:y}.

\input{build/floats/nb09_recovery}

The estimator recovers the truth: the mean across samples is \nbNineSeedMean{}\,pp against a planted
\nbNineTrueJump{}\,pp, a bias of \nbNineSeedBias{}\,pp. It is honest about its uncertainty: the
sample-to-sample standard deviation of the estimate is \nbNineSeedSd{}\,pp, and the nominal 90\,\%
interval covers the truth on \nbNineSeedCovN{} of \nbNineNSeeds{} samples --- \pc{\nbNineSeedCovPct},
essentially nominal. The \nbNineRdErr{}\,pp miss in this chapter's sample is therefore sampling
noise, and the interval was wide enough to say so.

That result matters for the chapter's argument, and it should be held next to what
\cref{sec:rd:wrong} is about to establish. \textbf{Conditional on the bandwidth, the interval is
well calibrated.} Every seed in \cref{fig:rd:recovery} was fitted at $h = \eur{\nbNineH}$. The
calibration study prices the sampling noise \emph{given} the window, because that is the only thing
it varied. It says nothing whatever about the window.

% =====================================================================================
\section{Point estimate versus posterior}
\label{sec:rd:compare}

\input{build/tables/nb09_compare}

\subsection*{They agree, and the chapter says so out loud}

\Cref{tab:rd:compare} puts both arms on the estimand of \cref{eq:rd:estimand}: the same jump, the
same \cref{as:rd:cont}, the same bandwidth. Two classical rows are shown rather than one, so that
the kernel cannot hide in the comparison --- CausalPy fits the window \emph{unweighted}, so the
exact classical twin of \cref{eq:rd:lik} is the uniform-kernel fit, while the triangular fit is the
canonical classical RD of \cref{sec:rd:classical}.

The three point estimates span \nbNineEstSpan{}\,pp. The classical triangular estimate is
\nbNineRd{}\,pp, the classical uniform estimate \nbNineUnif{}\,pp, the posterior mean
\nbNineBayes{}\,pp; all three intervals cover the planted \nbNineTrueJump{}\,pp; and the logistic
refit of \cref{sec:rd:bayes} agreed a fourth time at \nbNineLogitJump{}\,pp. \textbf{Bayes did not
move the number.}

Nor did it sharpen it. The posterior's half-width is \nbNineBayesHalf{}\,pp against
\nbNineUnifHalf{}\,pp for its exact classical twin --- a width ratio of \nbNineWidthRatio{}. With
\nbNineNWindow{} customers in the window and the effectively-flat priors of \cref{eq:rd:prior}, the
posterior \emph{is} the least-squares answer, re-expressed. This is the ordinary case in this book
and it has a name: the Bernstein--von Mises phenomenon \citep{gelman2013}. With a well-specified
likelihood and enough data, the posterior concentrates where the maximum-likelihood estimate does,
and the two paradigms return the same number with the same width. It is worth saying plainly,
because the marketing of Bayesian methods rarely does: \textbf{on this problem, the sampler bought
neither accuracy nor precision.}

That the widths agree to within \nbNineWidthRatio{} also disposes of the heteroskedasticity worry
raised in \cref{sec:rd:bayes}. The classical row \emph{does} correct for the Bernoulli variance and
the Bayesian row does not; the two intervals nevertheless land on top of one another, because with
retention probabilities in the middle of the unit interval, $p(1-p)$ varies little across the
window. A different application --- rare-event retention, with $p$ near zero on one side of the
cutoff --- would not get off so lightly, and there the constant-$\sigma$ likelihood would have to be
replaced rather than excused.

\subsection*{What the posterior did buy}

One thing, and it is the thing the CMO's question is made of. In the last column of
\cref{tab:rd:compare}, the classical entries are not ``small'' or ``hard to compute''. They are
\textbf{undefined}, and that is a literal statement rather than a rhetorical flourish. A confidence
interval is a statement about a procedure; a $p$-value ranks data against a null. Neither attaches a
probability to a hypothesis about $\tau_{\mathrm{RD}}$, and no amount of additional frequentist
machinery produces one.

The posterior does, immediately and at three thresholds that matter:
%
\begin{equation}
  \Prob(\tau_{\mathrm{RD}} > 0) = \nbNinePGtZero, \qquad
  \Prob(\tau_{\mathrm{RD}} > \nbNineBreakevenJump\,\text{pp}) = \nbNinePGtBe, \qquad
  \Prob(\tau_{\mathrm{RD}} > 10\,\text{pp}) = \nbNinePGtMaterial .
  \label{eq:rd:probs}
\end{equation}
%
The middle quantity is the perk's break-even jump, derived in \cref{sec:rd:euro}, and it is the only
number in this chapter that \cref{sec:rd:euro}'s decision rule consumes. The first says the effect is
almost certainly real. The third says that ``real'' and ``material'' are different questions with
different answers.

So the ledger, stated plainly. Bayes did not buy the point estimate; the classical one is as good.
It did not buy a tighter interval; the ratio is \nbNineWidthRatio. It did not buy a better
diagnostic; every test in \cref{sec:rd:valid} is classical. What it bought is a coherent way to turn
an estimate into a euro decision --- and, as \cref{sec:rd:wrong} insists, \textbf{it did not solve
the bandwidth problem either}. That problem is paradigm-independent, and a posterior conditioned on
$h = \eur{\nbNineH}$ is exactly as silent about the choice of $h$ as a confidence interval
conditioned on the same thing.

% =====================================================================================
\section{The decision, in euros}
\label{sec:rd:euro}

The perk costs \eur{\nbNinePerkCost} per member per year. A retained member is worth
\eur{\nbNineAnnualValue} a year. Roughly \nbNineNMargin{} members sit at the \eur{\nbNineCutoff}
margin --- the population \cref{eq:rd:estimand} is actually about, and the only one this chapter is
licensed to speak for.

The arithmetic per marginal member is one line. A jump of $\tau$ retains $\tau$ extra members per
member exposed, each worth \eur{\nbNineAnnualValue}, at a cost of \eur{\nbNinePerkCost}:
%
\begin{equation}
  \text{net value per member} \;=\; \tau \cdot v \;-\; k ,
  \qquad v = \nbNineAnnualValue, \quad k = \nbNinePerkCost ,
  \label{eq:rd:net}
\end{equation}
%
so the perk breaks even at $\tau^{\ast} = k / v = \nbNineBreakevenJump{}\,\text{pp}$. That is the
number \cref{eq:rd:probs}'s middle entry evaluates, and it is why the break-even was needed one
section early.

Propagating the posterior of $\tau_{\mathrm{RD}}$ through \cref{eq:rd:net} gives an expected net
value of \eur{\nbNineNetMember} per member and \eur{\nbNineTotalNet} across the \nbNineNMargin{}
marginal members, with a 90\,\% interval of $[\eur{\nbNineTotalLo},\ \eur{\nbNineTotalHi}]$ --- an
interval that \emph{straddles zero}. The decision rule is the book's standing convention:
%
\begin{equation}
  \Prob(\tau\, v > k)\ \begin{cases}
    > 0.9 & \Rightarrow\ \textsc{keep} \quad\text{(act without buying more evidence)},\\
    < 0.5 & \Rightarrow\ \textsc{kill} \quad\text{(more likely than not it does not pay)},\\
    \text{otherwise} & \Rightarrow\ \textsc{reconsider}.
  \end{cases}
  \label{eq:rd:rule}
\end{equation}
%
The posterior gives $\Prob(\text{perk pays}) = \nbNinePPays$ against the \nbNineActionBar{} bar. The
verdict is \textbf{\nbNineDecision} (\cref{fig:rd:euro}).

\input{build/floats/nb09_euro}

Note carefully \emph{what} makes it a \textsc{reconsider}, because this is where a decision memo
usually goes wrong. It is not the expected value, which is positive: \eur{\nbNineNetMember} per
member is real money at scale. It is the \textbf{spread}. The perk is more likely than not to pay,
and it is far from certain to. A firm that reads only the mean of \cref{fig:rd:euro}'s left panel
renews the programme; a firm that reads its width asks a better question.

\subsection*{The better question: what would settle it, and what would it cost?}

\textsc{reconsider} must not end as a shrug. There are exactly two levers, and the posterior prices
both.

\textbf{Lever one: buy more data.} The posterior standard deviation of the perk's gross value is
\eur{\nbNineSdNow} per member, and it shrinks roughly as $1/\sqrt{n}$ in the number of customers
inside the window. One more cycle of data --- roughly doubling the \nbNineNWindow{} members in the
$\pm$\eur{\nbNineH} band --- would tighten the credible interval by about
\pc{\nbNineCriShrink} and lift $\Prob(\text{pays})$ from \nbNinePPays{} to \nbNinePPaysNext{},
\emph{if the posterior mean holds}. Still short of the bar. Clearing \nbNineActionBar{} outright at
the current cost requires a posterior standard deviation of \eur{\nbNineSdReq}, which is
$\approx \nbNineDataFactor\times$ the current window sample. That is not one more cycle. It is
years.

\textbf{Lever two: renegotiate the price.} The right-hand panel of \cref{fig:rd:euro} sweeps the
perk's cost and reads $\Prob(\text{pays})$ off the same posterior. At a free perk it is
\nbNinePPaysFree; at \eur{\nbNinePerkGridHi} it is \nbNinePPaysDear. The cost that the
\emph{conservative} effect --- the posterior's 5th percentile --- can still bear is
\eur{\nbNineBePerk} per member. A perk costing \eur{\nbNineBePerk} rather than
\eur{\nbNinePerkCost} would pay even if the true effect sat at the pessimistic end of everything
this chapter has estimated.

\textbf{So the recommendation writes itself, and it is not ``keep'' or ``kill''.} The evidence
cannot be improved fast enough to matter, and it does not need to be: \emph{renegotiating the perk's
cost is the stronger lever by a wide margin}. Take \eur{\nbNineBePerk} to the supplier, or find a
perk with the same perceived value and a lower marginal cost. That is a deliverable a CMO can act on
this quarter, and it is more useful than either a $p$-value or a verdict.

\begin{remark}[What RD cannot answer, and the design that can]
The question the loyalty team will ask next is \emph{where should the threshold be?} ---
\eur{80}? \eur{150}? Regression discontinuity \textbf{cannot answer it}, and no amount of care with
this dataset will change that. \Cref{eq:rd:estimand} is identified at $s = c$ and nowhere else: the
effect $\tau(s)$ for customers at other spend levels is not a quantity the design sees, and moving
the threshold is an extrapolation beyond the only point where the method has any authority. Two
designs do answer it. Run the programme with \emph{several} cutoffs across different cohorts or
markets, which identifies $\tau(s)$ at each of them; or --- the clean answer --- \textbf{randomize
the cutoff itself}, assigning customers a threshold drawn from a range, which turns $\tau(\cdot)$
into an experimentally identified function rather than a single identified point. The honest sentence
in the deck is: \emph{this analysis prices the perk at the current margin, and pricing a different
margin requires a different design.}
\end{remark}

% =====================================================================================
\section{What can go wrong}
\label{sec:rd:wrong}

\subsection*{The bandwidth is a choice, not an estimate}

Everything above was computed at $h = \eur{\nbNineH}$. That number was chosen. It was not estimated,
it was not derived, and no interval in this chapter prices the consequences of having chosen it
differently.

Sweep it. \Cref{tab:rd:bwsweep} runs the identical estimator on the identical data at
\nbNineBwNWindows{} bandwidths from \eur{\nbNineBwMin} to \eur{\nbNineBwMax}, and
\cref{fig:rd:bandwidth} plots the result. Three things happen, in ascending order of seriousness.

\input{build/tables/nb09_bwsweep}
\input{build/floats/nb09_bandwidth}

\begin{enumerate}[leftmargin=*, itemsep=.4em]
  \item \textbf{The point estimate moves.} It ranges from \nbNineBwEstLo{}\,pp at the narrowest
        window to \nbNineBwEstHi{}\,pp at the widest --- a span of \nbNineBwEstRange{}\,pp, which is
        \pc{\nbNineBwRangeOverHalf} of the \emph{half-width} of the reported interval. The drift has
        a direction and a cause: narrow windows are noisy, wide windows pull in customers far from
        the margin whose retention is governed by the confounder's slope, and the local-linear fit
        starts booking curvature as level. This is the bias--variance trade-off, and it is the
        reason a bandwidth exists at all.
  \item \textbf{The significance verdict itself flips.} The narrow windows' intervals cross zero:
        \textbf{\nbNineBwNInsignif{} of the \nbNineBwNWindows{} windows cannot exclude it}
        (\cref{tab:rd:bwsweep}, column five). The same data, the same estimator, the
        same assumptions --- and whether the perk has a demonstrable effect at all depends on a knob
        the analyst turned before anyone was looking. Had the analysis been written up at
        $h = \eur{\nbNineBwMin}$, the finding would have been ``no significant effect'' and the perk
        would have been killed.
  \item \textbf{The union of the intervals is enormous.} Taking the lowest lower bound and the
        highest upper bound across the sweep gives $[\nbNineBwUnionLo, \nbNineBwUnionHi]$\,pp ---
        \nbNineBwUnionWidth{}\,pp wide, against the \nbNineRdWidth{}\,pp of the single reported
        interval, a factor of \nbNineBwUnionOverWidth. \textbf{That is the width the reported
        interval is not showing you.} The shaded band in \cref{fig:rd:bandwidth} is that union, and
        nothing inside the reported interval knows it exists.
\end{enumerate}

Now state the general principle, because it is much bigger than this chapter.

\begin{quote}
\textbf{An interval computed conditional on a tuning parameter prices sampling noise \emph{given}
that parameter. It does not price the choice of the parameter.} Plug in a hand-chosen $h$, report
the interval the estimator returns, and the uncertainty attributable to $h$ is silently set to zero
--- not estimated as small, but omitted from the arithmetic entirely.
\end{quote}

That is a property of \emph{plugging in}, and it has nothing to do with bandwidths in particular. The
marketing-mix chapter of this book reports a confident credible interval on ROAS --- computed with
the adstock decay and the saturation curvature fixed by hand. Sweep those two and the ROAS interval
moves in exactly the way \cref{fig:rd:bandwidth} moves, and for exactly the same reason. The same
holds of a matching estimator's caliper, a synthetic control's donor-pool filter, and a
regularization penalty picked by eye. \Citet{leamer1983} named the disease forty years ago: the
reported specification is one draw from a cloud of specifications the analyst ran, and the interval
describes the draw rather than the cloud.

\textbf{Neither paradigm rescues this, and it is important to be exact about why.} The classical
interval is conditional on $h$. The posterior is conditional on $h$. \emph{Bayes did not solve the
bandwidth problem} --- putting a prior on $\tau$ says nothing about $h$, because $h$ never appeared
in \crefrange{eq:rd:lik}{eq:rd:prior} as a parameter at all. It was baked into the data the
likelihood saw. A Bayesian who genuinely wanted to price it would have to make $h$ a parameter and
integrate over it, or average across models indexed by $h$ --- which is a different model from the
one this chapter fitted, and one whose prior over $h$ would then be doing real work for the first
time in this book.

What honest practice looks like, given all that:

\begin{itemize}[leftmargin=*, itemsep=.3em]
  \item \textbf{Report the sweep, not the point.} \Cref{tab:rd:bwsweep} belongs in the appendix of
        every RD deck ever written. A reader who sees only ``\nbNineRd{}\,pp
        $[\nbNineRdLo, \nbNineRdHi]$'' has been told less than the analysis knows.
  \item \textbf{Use a data-driven bandwidth, and then still sweep.} \Citet{imbens2012} derive the
        $h$ that minimises mean-squared error, and \citet{calonico2014} show that the resulting
        estimator is biased at exactly the rate that invalidates the naive interval --- so they
        supply a bias-corrected estimator with a robust interval that remains valid at the
        MSE-optimal $h$. That is the state of the art, it is what \code{rdrobust} implements, and
        it is what a production RD analysis should use. It narrows the choice; it does not abolish
        it, because the optimal $h$ is itself estimated and the sweep still moves.
  \item \textbf{Pre-register the bandwidth.} The cheapest fix in the list, and the most effective:
        choose $h$ before looking at the outcome, and the researcher degree of freedom disappears
        by construction.
\end{itemize}

\subsection*{High-order global polynomials: the remedy that is worse}

There is an old temptation, when the local-linear fit looks wobbly, to widen the window to the whole
sample and control for curvature with a global polynomial of order four, five or six. It fits
beautifully. It is also, as \citet{gelman2019} demonstrate, a mistake of a specific and serious
kind. High-order polynomials give enormous weight to observations \emph{far} from the cutoff ---
the whales, who have nothing to say about \cref{eq:rd:estimand} --- and the resulting estimator has
noisy, wildly varying implicit weights near the boundary. Runge's phenomenon does the rest: the
fitted curve oscillates near the edges of its range, and the cutoff \emph{is} an edge for each of
the two fits. The estimand it targets is a jump at $c$; what it delivers is a difference of two
extrapolated polynomial endpoints. \emph{Do not do it.} Narrow the window and keep the polynomial
low-order; the quadratic check of \cref{sec:rd:valid} is as far up as this chapter goes, and it went
there only to confirm that the linear fit did not need help.

\subsection*{Heaping, and the donut}

Real running variables are not smooth draws from a normal distribution. Spend heaps at round numbers
--- \eur{100} exactly, \eur{99.99}, \eur{50} --- because prices and promotions heap there, and a
threshold placed \emph{at} a heap point mixes the customers who landed on \eur{100} for a dozen
unrelated reasons into the comparison at precisely the place where the estimator is most sensitive.
The standard remedy is a \textbf{donut RD}: discard the observations within $\pm\epsilon$ of the
cutoff and refit, so that the two local-linear fits are extrapolated a short distance to the
threshold from data that are not contaminated by the heap. The simulator here has no heaping, so
there is nothing to demonstrate --- but on real spend data the first thing to plot is a histogram at
one-euro resolution, and the second is the donut.

\subsection*{The rest of the list}

\emph{The effect is local, and the estimand says so.} \Cref{eq:rd:estimand} is about the customer at
\eur{\nbNineCutoff}. Quoting \nbNineRd{}\,pp as ``the effect of Gold'' and multiplying it by the
entire Gold membership is a category error with a budget attached; the correct multiplicand is the
\nbNineNMargin{} members at the margin, which is what \cref{sec:rd:euro} used.

\emph{RD is expensive in data.} It discarded \pc{\nbNineRdDroppedPct} of the customer base
(\cref{tab:rd:naive}), and every diagnostic in \cref{sec:rd:valid} was computed on what remained. A
firm with a threshold and only a few hundred customers near it does not have an RD design; it has a
scatter plot.

\emph{And the cutoff must be sharp in practice, not merely on paper.} If Gold is granted at year end
on a spend figure that is later restated, if service agents can grant it as a goodwill gesture, or
if the tier persists for a year after the customer's spend falls back --- then $D$ is not
$\mathbf 1[S \ge c]$, \cref{eq:rd:running} is fiction, and the design is fuzzy at best. Read the
programme rules before writing the code.

% =====================================================================================
\section{Takeaways}
\label{sec:rd:take}

\begin{enumerate}[leftmargin=*, itemsep=.45em]
  \item \textbf{A tight interval around the wrong number is the most dangerous object in
        analytics.} The Gold-vs-rest gap is \nbNineNaive{}\,pp with an interval of
        $[\nbNineNaiveLo, \nbNineNaiveHi]$, and the truth is \nbNineNaiveVerdict{} it
        (\cref{tab:rd:naive}). The interval is not lying: it prices sampling noise correctly, and
        \emph{nothing in a confidence interval knows about confounding}. It will pass every review.
        Confounding is not a variance problem and no volume of data fixes it.

  \item \textbf{A sharp threshold is an experiment nobody ran.} Just below and just above
        \eur{\nbNineCutoff}, customers are essentially identical, and the only thing that changes
        discontinuously is the perk. Reading only the vertical break at the cutoff
        (\cref{fig:rd:fit}) cancels the confounder that the naive comparison booked as effect ---
        \nbNineRd{}\,pp $[\nbNineRdLo, \nbNineRdHi]$, which contains the planted
        \nbNineTrueJump{}\,pp.

  \item \textbf{Comparability is bought with data, and the price is precision.} The estimate rests
        on \nbNineRdN{} customers; \nbNineRdDropped{} of \nbNineN{} were thrown away, and the
        interval is \nbNineWidthVsNaive$\times$ as wide as the naive one. \emph{Precision you did
        not earn is not a virtue}, and an analyst who optimises for narrow intervals will pick the
        wrong estimator every time.

  \item \textbf{RD's diagnostics are unusually good --- use all of them.} The density is smooth at
        the cutoff ($z = \nbNineMccZ$), pre-determined covariates do not jump
        ($t = \nbNineCovTenureT$, $t = \nbNineCovOrdersT$), \nbNinePlaceboN{} placebo cutoffs find
        nothing (max $|t| = \nbNinePlaceboMaxT$), and the quadratic refit moves the answer by
        \nbNinePolyDiff{}\,pp. Each is a test of a \emph{named} assumption
        (\cref{tab:rd:assumptions}) that could have failed and did not. The central assumption ---
        continuity of the potential-outcome means, \cref{as:rd:cont} --- \emph{cannot} be tested
        directly. Only its implications can, and saying so is part of the analysis.

  \item \textbf{The bandwidth is a choice, not an estimate, and no interval in this chapter prices
        it.} Across \nbNineBwNWindows{} windows the estimate ranges \nbNineBwEstLo{} to
        \nbNineBwEstHi{}\,pp and \nbNineBwNInsignif{} of them cannot exclude zero --- the
        significance verdict itself flips. The union of the intervals is
        $[\nbNineBwUnionLo, \nbNineBwUnionHi]$\,pp, \nbNineBwUnionOverWidth$\times$ the reported
        width. This is a general property of \emph{plugging in} a tuning parameter, not a quirk of
        RD: it is the marketing-mix model's hand-fixed adstock, the matching estimator's caliper,
        the penalty chosen by eye. Report the sweep; use \citeauthor{calonico2014}'s bias-corrected
        estimator at an \citeauthor{imbens2012} bandwidth; pre-register $h$ where you can.

  \item \textbf{Bayes did not move the number, did not sharpen it, and did not solve the bandwidth
        problem.} The point estimates span \nbNineEstSpan{}\,pp and the posterior is
        \nbNineWidthRatio$\times$ the width of its exact classical twin (\cref{tab:rd:compare}) ---
        Bernstein--von Mises, exactly as advertised, because with \nbNineNWindow{} customers and
        flat priors the posterior \emph{is} the least-squares answer re-expressed. And a posterior
        conditioned on $h = \eur{\nbNineH}$ is as silent about the choice of $h$ as a confidence
        interval conditioned on it. What the posterior bought is the one object the decision
        consumes: $\Prob(\tau_{\mathrm{RD}} > \nbNineBreakevenJump\,\text{pp}) = \nbNinePGtBe$, which
        has no classical counterpart at any level of effort.

  \item \textbf{The decision is about the price, not the evidence.} $\Prob(\text{perk pays})$ is
        \nbNinePPays{} against a \nbNineActionBar{} bar: \textsc{\nbNineDecision}. Buying the
        evidence to clear that bar would take $\approx \nbNineDataFactor\times$ the current sample
        --- one more cycle only reaches \nbNinePPaysNext. But the perk pays comfortably at
        \eur{\nbNineBePerk} per member instead of \eur{\nbNinePerkCost}. \emph{Renegotiating the
        cost is the stronger lever}, and the analysis that says so is more useful than one that says
        yes or no.

  \item \textbf{RD cannot tell you where to put the threshold.} $\tau(s)$ is identified at $s = c$
        and nowhere else. ``Should Gold start at \eur{80}?'' is a question about a point the design
        never sees, and answering it from this data is extrapolation dressed as inference. The
        design answer is to run several cutoffs --- or to randomize the cutoff itself.
\end{enumerate}

\notebooknote{%
  The analysis of \crefrange{sec:rd:dgm}{sec:rd:wrong} is
  \code{notebooks/09\_threshold\_perk\_rdd.ipynb}, which runs in the legacy environment
  (\code{pymc < 6}, kernel \code{cmp-legacy}) because it uses CausalPy's
  \code{RegressionDiscontinuity}. It has a fast teaching mode (\code{CMP\_FAST=1}, short chains,
  fewer recovery seeds) and a full mode (\code{CMP\_FAST=0}); every number in this chapter comes
  from the full run, emitted by \code{cmp.report} and injected here as a macro. The simulator is
  \code{cmp.dgp.rdd\_perk}, the classical estimator is \code{cmp.classical.rd\_local\_linear}, the
  density check is \code{cmp.estimators.mccrary\_density}, and the seed is fixed at
  \code{SEED = \nbNineSeed}. For production work, \citeauthor{calonico2014}'s \code{rdrobust} is the
  reference implementation, and \citet{cattaneo2020} is the practical introduction to read next.}
